% Beamer presentation for SIR Project
\documentclass{beamer}

% Theme
\usetheme{Madrid}
\usecolortheme{seahorse}

% Packages
\usepackage[utf8]{inputenc}
\usepackage{graphicx}
\usepackage{amsmath}
\usepackage{amssymb}
\usepackage{babel}

\title[Modelo SIR y MCMC]{Modelo SIR como Cadena de Markov aplicado al dengue en Cali 2024}
\author{Elián Steven Bejarano González
\\ Cadenas de Markov y Aplicaciones}
\institute{Universidad Nacional de Colombia}
\date{2025}

\begin{document}

% Slide 1 --------------------------------------------------------------------
\begin{frame}
  \titlepage
\end{frame}

% Slide 2 --------------------------------------------------------------------
\begin{frame}{Motivación: Dengue en Cali 2024}
  \begin{itemize}
    \item Incremento histórico de casos confirmados.
    \item Datos obtenidos de microdatos SIVIGILA.
    \item Caso ideal para estudiar modelos epidémicos estocásticos.
  \end{itemize}
  \begin{center}
    \includegraphics[width=0.75\textwidth]{fig_incidencia_cali.png}
  \end{center}
\end{frame}

% Slide 3 --------------------------------------------------------------------
\begin{frame}{Modelo SIR: Estructura general}
  \begin{columns}
    \column{0.55\textwidth}
      \begin{itemize}
        \item Estado del sistema:
        \[
          X_t = (S_t, I_t, R_t)
        \]
        \item Propiedad de Markov.
        \item Transiciones:
        \[
          \Delta I_t \sim \text{Binomial}(S_t, \beta I_t)
        \]
        \[
          \Delta R_t \sim \text{Binomial}(I_t, \gamma)
        \]
      \end{itemize}
    \column{0.45\textwidth}
      \begin{center}
      \includegraphics[width=0.9\textwidth]{fig_flechas_sir.png}
      \end{center}
  \end{columns}
\end{frame}

% Slide 4 --------------------------------------------------------------------
\begin{frame}{Implementación Bayesiana con JAGS}
  \begin{itemize}
    \item Ajuste completo mediante MCMC.
    \item 4 cadenas independientes.
    \item 80,000 iteraciones de adaptación.
    \item Posteriores para $S_0$, $\beta$, $\tau_C$.
  \end{itemize}
  \begin{center}
    \includegraphics[width=0.75\textwidth]{fig_modelo_jags.png}
  \end{center}
\end{frame}

% Slide 5 --------------------------------------------------------------------
\begin{frame}{Diagnóstico MCMC}
  \begin{center}
    \includegraphics[width=0.9\textwidth]{fig_traceplots_parametros.png}
  \end{center}
  \begin{itemize}
    \item Trazas no unimodales.
    \item Saltos entre modos en $S_0$ y $\beta$.
    \item Evidencia de multimodalidad: el modelo SIR null es insuficiente.
  \end{itemize}
\end{frame}

% Slide 6 --------------------------------------------------------------------
\begin{frame}{Resultados: Observados vs Modelo}
  \begin{center}
    \includegraphics[width=0.9\textwidth]{fig_casos_vs_modelo.png}
  \end{center}
  \begin{itemize}
    \item El modelo reproduce la forma general del brote.
    \item La magnitud predicha es irreal: sobreestima enormemente.
  \end{itemize}
\end{frame}

% Slide 7 --------------------------------------------------------------------
\begin{frame}{Trayectorias S(t) y R(t)}
  \begin{columns}
    \column{0.5\textwidth}
      \includegraphics[width=0.95\textwidth]{fig_susceptibles.png}
    \column{0.5\textwidth}
      \includegraphics[width=0.95\textwidth]{fig_recuperados.png}
  \end{columns}
  \begin{itemize}
    \item Caída brusca en susceptibles.
    \item Inmunidad casi total → irreal en contexto urbano real.
  \end{itemize}
\end{frame}

% Slide 8 --------------------------------------------------------------------
\begin{frame}{Discusión}
  \begin{itemize}
    \item El modelo SIR null no captura estacionalidad.
    \item La multimodalidad posterior indica incompatibilidad estructural.
    \item Motiva modelos con $\beta(t)$ o clima.
  \end{itemize}
\end{frame}

% Slide 9 --------------------------------------------------------------------
\begin{frame}{Conclusiones}
  \begin{itemize}
    \item El SIR null es útil pedagógicamente, pero no para dengue real.
    \item MCMC revela claramente las limitaciones estructurales.
    \item Extensiones: clima, modelos vectoriales, $\beta(t)$.
  \end{itemize}
\end{frame}

% Slide 10 --------------------------------------------------------------------
\begin{frame}{Gracias}
  \begin{center}
    \Large Gracias por su atención.
    \\[0.5cm]
    \small Código completo en GitHub https://github.com/elbejaranog-runner/Modelado-estoc-stico-del-dengue-en-Cali-mediante-un-modelo-SIR-y-Cadenas-de-Markov).
  \end{center}
\end{frame}

\end{document}
